Linear Allocator (линейный аллокатор) "--- это самый простой вид аллокаторов. Его идея заключается в том, чтобы сохранить указатель на начало блока памяти, за который отвечает аллокатор, а также использовать другой указатель, который необходимо будет перемещать каждый раз, когда завершено очередное выделение памяти. Этот аллокатор позволяет почти полностью избежать внутренней фрагментации, поскольку элементы вставляются последовательно, и единственная фрагментация между ними "--- выравнивание \cite{github-allocators}. Основной его минус заключается в том, что мы не можем освободить память, выделенную под конкретный объект, не обнулив весь аллокатор.

Мы сохраняем указатель на начало памяти (\texttt{start}), а также храним два числа, которые содержат информацию об общем (\texttt{end}) и используемом (\texttt{used}) размерах памяти. При запросе выделения блока памяти размера \texttt{k} аллокатором будут выполнены следующие действия:

\begin{enumerate}
	\item проверка, достаточно ли памяти для выделения ($\texttt{used} + \texttt{k} \le \texttt{end}$);
    \item сохранение текущего указателя \texttt{used}, который в дальнейшем будет отдан пользователю как указатель на блок выделенной памяти;
	\item смещение указателя used на величину, равную объёму выделенного блока памяти ($\texttt{used} = \texttt{used} + \texttt{k}$).
\end{enumerate}

Для очистки памяти необходимо установить \texttt{used} в $0$ \cite{allocators}.

Применения линейного аллокатора:

\begin{itemize}
	\item В игровых движках:
	\begin{itemize}
		\item Выделение памяти под временные данные, например, для рендеринга кадра. В конце кадра вся память освобождается, и аллокатор сбрасывается.
		\item В процессе расчётов, связанных с физикой объекта (столкновение, движение), для моделирования системы частиц.
		\item При обработке звука для создания временных буферов \cite{game++}.
	\end{itemize}
	\item При обработке структур данных (JSON, деревья): память под узлы выделяется последовательно, после окончания обработки она вся освобождается.
\end{itemize}
